\documentclass[sigplan]{acmart}

\setcopyright{rightsretained}
\copyrightyear{2024}
\acmYear{2024}
\acmConference[Incorrectness '24]{Formal Methods for Incorrectness}{January 2024}{London, UK}
\acmBooktitle{Incorrectness '24: Formal Methods for Incorrectness}
\acmDOI{}
\acmISBN{}
\expandafter\def\csname @copyrightpermission\endcsname{\raisebox{-2ex}{\includegraphics[width=.2\linewidth]{cc-by}} \parbox{.7\linewidth}{\raggedright This work is licensed under a Creative Commons Attribution 4.0 International License.}}
\expandafter\def\csname @copyrightowner\endcsname{Roblox.}
\newcommand{\infer}[2]{\frac{\displaystyle\begin{array}{@{}l@{}}#1\end{array}}{\displaystyle#2}}
\newcommand{\LOCAL}{\mathop{\mathsf{local}}}
\newcommand{\FUNCTION}{\mathop{\mathsf{function}}}
\newcommand{\IF}{\mathop{\mathsf{if}}}
\newcommand{\THEN}{\mathbin{\mathsf{then}}}
\newcommand{\ELSE}{\mathbin{\mathsf{else}}}
\newcommand{\END}{\mathop{\mathsf{end}}}
\newcommand{\NEVER}{\mathsf{never}}
\newcommand{\ERROR}{\mathsf{error}}
\newcommand{\UNKNOWN}{\mathsf{unknown}}
\newcommand{\STRING}{\mathsf{string}}
\newcommand{\NUMBER}{\mathsf{number}}
\newcommand{\MATH}{\mathsf{math}}
\newcommand{\ABS}{\mathsf{abs}}
\newcommand{\LOWER}{\mathsf{lower}}
\newcommand{\fun}{\mathbin{\rightarrow}}

\begin{document}

\title{Towards an Unsound But Complete Type System}
\subtitle{Work In Progress on New Non-Strict Mode for Luau}

\author{Lily Brown}
\author{Andy Friesen}
\author{Alan Jeffrey}
\author{Vighnesh Vijay}
\affiliation{
  \institution{Roblox}
  \city{San Mateo}
  \state{CA}
  \country{USA}
}

\begin{abstract}
In HATRA 2021, we presented \emph{The Goals Of The Luau Type System},
describing the human factors of a type system for a language with a
heterogeneous developer community.  One of the goals was the design of
type system for bug detection, where we have high confidence that type
errors isolate genuine software defects, and that false positives are
minimized. Such a type system is, by necessity, unsound, but we can ask
instead that it is complete. This paper presents a work-in-progress report
on the design and implementation of the new non-strict mode for Luau.
\end{abstract}

\maketitle

\section{Introduction}

Luau~\cite{Luau} is the scripting language used by the
Roblox~\cite{Roblox} platform for shared immersive experiences.  Luau extends
the Lua~\cite{Lua} language, notably by providing type-driven tooling
such as autocomplete and API documentation (as well as traditional type
error reporting). Roblox has hundreds of millions of users, and
millions of creators, ranging from children learning to program for
the first time to professional development studios.

In HATRA 2021, we presented a position paper on the \emph{Goals Of The Luau Type
System}~\cite{BFJ21:GoalsLuau}, describing the human factors issues
with designing a type system for a language with a heterogeneous
developer community. The design flows from the needs of the different
communities: beginners are focused on immediate goals (``the stairs
should light up when a player walks on them'') and less on the code
quality concerns of more experienced developers; for all users
type-driven tooling is important for productivity. These needs result in a design with two modes:
\begin{itemize}
\item \emph{non-strict mode}, aimed at non-professionals, focused on
  minimizing false positives (that is, in non-strict mode, any program with a type error has a defect), and
\item \emph{strict mode}, aimed at professionals, focused on
  minimizing false negatives (that is, in strict mode, any program with a defect has a type error).
\end{itemize}
The focus of this extended abstract is the design of non-strict mode: what constitutes
a defect, and how can we design a complete type system for detecting them.
(The words ``sound'' and ``complete'' in this sense arte far from ideal,
but ``sound type system'' has a well-established meaning, and ``complete''
is well-established as the dual of ``sound'', so here we are.)

The closest work to ours is success typing~\cite{SuccessTyping}, used
in Erlang Dialyzer~\cite{Dialyzer}.  The new feature of our work is
that strict and non-strict mode have to intertact, whereas Dialyzer only has
the equivalent of non-strict mode.

New non-strict mode is specified in a Luau Request For
Comment~\cite{NewNonStrictRFC}, and is currently being implemented.
We expect it (and other new typechecking features) to be available in
2024. This extended abstract is based on the RFC, but written in
``greek letters and horizontal lines'' rather than ``monspaced text''
for mat.

\section{Code defects}

The main goal of non-strict mode is to identify defects, but this requires
identifying what constitutes a defect,. Run-time errors are an obvious defect:
\begin{verbatim}
  local hi = "hi"
  print(math.abs(hi))
\end{verbatim}
but we also want to catch common mistakes such as mis-spelling a property name,
even though Luau returns \verb|nil| for missing property accesses.
For this reason, we consider a larger class of defects:
\begin{itemize}
\item run-time errors,
\item expressions guaranteed to be \verb|nil|, and
\item writing to a table property that is never read.
\end{itemize}

\subsection{Run-time errors}

Run-time errors occur due to run-time type mismatches (such as \verb|5("hi")|)
or buit-in functios (such as \verb|math.abs("hi")|).
Detecting run-time errors is undecidable, for example:
\begin{verbatim}
  if cond() then
    math.abs(“hi”)
  end
\end{verbatim}
In general, we cannot be sure that this code produces a run-time
error, but we do know that if \verb|math.abs("hi")| is executed, it
will produce an error, so we consider this to be a defect.

\subsection{Expressions guaranteed to be nil}

Luau tables do not error when a missing property is accessed (though embeddings may). So
\begin{verbatim}
  local t = { Foo = 5 }
  local x = t.Fop
\end{verbatim}
won’t produce a run-time error, but is more likely than not a
programmer error. In this case, if the programmer intent was to
initialize \verb|x| as \verb|nil|, they could have initialized
\verb|x = nil|.  For this reason, we consider it a code defect to use
an expression that the type system infers is of type \verb|nil|, other
than the \verb|nil| constant itself.

\subsection{Writing properties that are never read}

There is a matching problem with misspelling properties when writing. For example
\begin{verbatim}
  function f()
    local t = {}
    t.Foo = 5
    t.Fop = 7
    print(t.Foo)
  end
\end{verbatim}
won’t produce a run-time error, but is more likely than not a
programmer error, since \verb|t.Fop| is written but never read. We can use
read-only and write-only table properties types for this, and make it an
error to create a write-only property.

We have to be careful about this though, because if \verb|f| ended
with \verb|return t|, then it would be a perfectly sensible function
with type \verb|() -> { Foo: number, Fop: number }|. The only way to detect
that \verb|Fop| was never read would be whole-program analysis, which is
prohibitively expensive.

\section{New Non-strict error reporting}

The difficult part of non-strict mode error-reporting is detecting
guaranteed run-time errors. We do this using an error-reporting
pass that generates a type context such that if any of the $x : T$ in
the type context are satisfied, then the program is guaranteed to
produce a type error.

For example in the program
\begin{verbatim}
  function h(x, y)
    math.abs(x)
    string.lower(y)
  end
\end{verbatim}
an error is reported when \verb|x| isn’t a \verb|number|, or \verb|y| isn’t a \verb|string|, so the generated context is
\begin{verbatim}
  x : ~number
  y : ~string
\end{verbatim}
(\verb|~T| is Luau's concrete syntax for type negation.)
In the function:
\begin{verbatim}
  function f(x)
    math.abs(x)
    string.lower(x)
  end
\end{verbatim}
an error is reported when \verb|x| isn’t a \verb|number| or isn’t a \verb|string|, so the constraint set is
\begin{verbatim}
  x : ~number | ~string
\end{verbatim}
(\verb"T|U" is Luau's concrete syntax for type negation.)
Since \verb"~number | ~string" is equivalent to the type \verb|unknown| (which contains every value),
non-strict mode can report a warning, since calling the function is guaranteed to throw a run-time error.
In contrast:
\begin{verbatim}
  function g(x)
    if cond() then
      math.abs(x)
    else
      string.lower(x)
    end
  end
\end{verbatim}
generates context
\begin{verbatim}
  x : ~number & ~string
\end{verbatim}
(\verb|T&U| is Luau's concrete syntax for type intersection.)
Since \verb|~number & ~string| is not equivalent to \verb|unknown|, non-strict mode reports no warning.

\begin{figure*}
  \[\begin{array}{c}
    \infer{
      \Gamma \vdash M : \NEVER \dashv \Delta_1 \\
      \Gamma, x : T \vdash B \dashv \Delta_2, x : U \\
      \mbox{(warn if $\UNKNOWN <: U$)}
    }{
      \Gamma \vdash (\LOCAL x : T = M; B) \dashv (\Delta_1 \cup \Delta_2)
    }
    \quad
    \infer{
      \Gamma \vdash M : \NEVER \dashv \Delta_1 \\
      \Gamma \vdash B \dashv \Delta_2 \\
      \Gamma \vdash C \dashv \Delta_3 
    }{
      \Gamma \vdash (\IF M \THEN B \ELSE C \END) \dashv (\Delta_1 \cup (\Delta_2 \cap \Delta_3))
    }
  \\[\bigskipamount]
    \infer{}{
      \Gamma \vdash x : T \dashv (x : T)
    }
    \quad
    \infer{
      \mbox{(warn if $k:T$)}
    }{
      \Gamma \vdash k : T \dashv \emptyset
    }
    \quad
    \infer{
      \Gamma, x:S \vdash B \dashv \Delta, x:U \\
      \mbox{(warn if $\UNKNOWN <: U$)}\\
      \mbox{(warn if ${\FUNCTION} <: T$)}
    }{
      \Gamma \vdash (\FUNCTION (x : S) B \END) : T \dashv \Delta
    }
    \quad
    \infer{
      \Gamma \vdash M : (S \fun \ERROR) \\
      \Gamma \vdash M : \neg{\FUNCTION} \dashv \Delta_1 \\
      \Gamma \vdash N : S \dashv \Delta_2  \\
      \mbox{(warn if $\Gamma \vdash M : (\UNKNOWN \fun T)$)}
    }{
      \Gamma \vdash M(N) : T \dashv \Delta_1 \cup \Delta_2
    }
  \end{array}\]
  \caption{Type context generation for blocks ($\Gamma \vdash B \dashv \Delta$) and expressions ($\Gamma \vdash M:T \dashv \Delta$)}
  \label{fig:ctxtgen}
\end{figure*}

\begin{figure*}
  \[
    \infer{
      \infer{}{\Gamma \vdash \MATH.\ABS : (\neg\NUMBER \fun \ERROR)} \quad
      \infer{}{\Gamma \vdash \MATH.\ABS : \neg{\FUNCTION} \dashv \emptyset} \quad
      \infer{
        \Gamma \vdash \STRING.\LOWER : (\neg\STRING \fun \ERROR) \\
        \Gamma \vdash \STRING.\LOWER : \neg{\FUNCTION} \dashv \emptyset \\
        \Gamma \vdash x : \neg{\STRING} \dashv (x : \neg\STRING) \\
        \mbox{(warn since $\Gamma \vdash \STRING.\LOWER : \UNKNOWN \fun \neg\NUMBER$)}
      }{\Gamma \vdash \STRING.\LOWER(x) : \neg{\NUMBER} \dashv (x : \neg\STRING)}
    }{\Gamma \vdash (\MATH.\ABS(\STRING.\LOWER(x)) : \NEVER \dashv (x : \neg\STRING)}
  \]
  \caption{Example warning}
  \label{fig:example}
\end{figure*}


In Figure~ref{fig:ctxtgen} we provide some of the inference rules for
conbtext generation, and the warnings that context generation produces.
These are run after type inference, so theay can assume that all code is fully typed.

These rules generaralize type intersection and union to type contexts,
write $\emptyset$ for the everywhere-$\NEVER$ context, and write $x:T$
for the context that is everywhere $\NEVER$ except for $x$ where it
maps to $T$.

\begin{conjecture}
If $\Gamma \vdash M : T \dashv \Delta, x:U$ and $\sigma$ is a closing
substitution where $\sigma(x) : U$, and $M[\sigma] \rightarrow^* v$
then $v : T$.
\end{conjecture}

\begin{corollary}
If $\Gamma \vdash M : \NEVER \dashv \Delta, x:\UNKNOWN$ and $\sigma$ is a closing
substitution and $M[\sigma]$ does not terminate successfully.
\end{corollary}

\section{Checked functions}

TODO

\section{Future work}

TODO

\bibliographystyle{ACM-Reference-Format} \bibliography{bibliography}

\end{document}
